\documentclass[11pt,a4paper]{article}

\usepackage[a4paper,margin=1in]{geometry}
\usepackage{fontspec}
\IfFontExistsTF{Times New Roman}
  {\setmainfont{Times New Roman}}
  {\IfFontExistsTF{TeX Gyre Termes}{\setmainfont{TeX Gyre Termes}}{\setmainfont{Times}}}
\usepackage{newtxmath}
\usepackage{amsmath}
\usepackage{booktabs}
\usepackage{longtable}
\usepackage{array}
\usepackage{enumitem}
\usepackage{xcolor}
\usepackage{listings}
\usepackage[round,authoryear]{natbib}
\usepackage[hidelinks]{hyperref}
\usepackage{xurl}
\usepackage{fancyhdr}
\usepackage{microtype}

\setlength{\parindent}{0pt}
\setlength{\parskip}{0.65em}
\setlist[itemize]{leftmargin=1.4em,itemsep=0.2em,topsep=0.2em}
\setlist[enumerate]{leftmargin=1.6em,itemsep=0.2em,topsep=0.2em}

\pagestyle{fancy}
\fancyhf{}
\lhead{\small Ticket 04 Study Report}
\rhead{\small PSOR/LCP Core}
\cfoot{\thepage}
\renewcommand{\headrulewidth}{0.3pt}

\lstset{
  basicstyle=\ttfamily\small,
  breaklines=true,
  frame=single,
  columns=fullflexible,
  keepspaces=true,
  rulecolor=\color{black!30},
  xleftmargin=0.5em,
  xrightmargin=0.5em
}

\newcommand{\code}[1]{\texttt{#1}}
\newcommand{\file}[1]{\path{#1}}
\newcolumntype{L}[1]{>{\raggedright\arraybackslash}p{#1}}
\newcolumntype{P}[1]{>{\raggedright\arraybackslash\ttfamily\footnotesize}p{#1}}

\title{\textbf{Ticket 04 Study Report: PSOR/LCP Core for American Option Obstacle}\\
\large SURF2026 American Option Risk Surfaces}
\author{Codex-assisted implementation report}
\date{June 16, 2026}

\begin{document}
\maketitle

\begin{abstract}
Ticket 04 adds the first American-option numerical core: projected successive
over-relaxation, or PSOR, for the linear complementarity problem created by the
American payoff obstacle. It also adds a basic Crank-Nicolson/PSOR time-marching
solver for American puts and dividend-paying American calls. This is a reusable
solver layer, not the full American put validation study. The report explains the
finance and numerical ideas slowly for beginner researchers who know Python and
data work, but are still learning variational inequalities, LCPs, PSOR, and
computational finance.
\end{abstract}

\tableofcontents
\newpage

\section{Role of Ticket 04 in the Whole Solver Project}

Tickets 01 through 03 built the analytic and European finite-difference foundation. Ticket 01
implemented payoff and European Black-Scholes closed-form utilities. Ticket 02 implemented the
uniform grids, finite-difference spatial operator, and boundary-value helpers. Ticket 03
implemented European Crank-Nicolson validation against closed-form Black-Scholes prices.

Ticket 04 is the first ticket that handles the American option obstacle. American options can be
exercised before maturity, so their value must not fall below immediate exercise payoff. This
turns the numerical problem from a plain linear system into a constrained problem. Ticket 04 adds
the reusable projected SOR logic needed to enforce that constraint.

This ticket is deliberately limited. It implements the PSOR/LCP core and a basic American
CN/PSOR time-marching solver. It does not perform the full American put validation study. Ticket
05 will use this core to build validation tables, representative price checks, and deeper
obstacle analysis.

\section{Theory and Methodology Connection}

The source basis for Ticket 04 is internal and project-specific:
\begin{itemize}
  \item \file{reports/01_literature/literature_map.md};
  \item \file{reports/02_math/formulation_note.md};
  \item \file{reports/03_solver/solver_validation_plan.md};
  \item \file{reports/03_solver/solver_implementation_tickets.md};
  \item the Ticket 01, Ticket 02, and Ticket 03 reports and code.
\end{itemize}

The literature map identifies variational inequalities, linear complementarity problems, and PSOR
as central concepts in classical finite-difference American option pricing. This reference pass
links those claims to verified sources: Brennan-Schwartz finite-difference American put valuation
\citep{brennan_schwartz_1977}, general LCP methodology \citep{cottle_pang_stone_2009}, and the
Black-Scholes/Merton option-pricing framework \citep{black_scholes_1973,merton_1973}. The
internal literature map, formulation note, and solver validation plan remain the implementation
source basis.

The methodology connection is stronger in Ticket 04 than in Tickets 01 through 03 because PSOR is
where the American exercise right enters the numerical algorithm. Ticket 01 supplied payoff and
European closed-form reference functions. Ticket 02 supplied the spatial grid and Black-Scholes
operator. Ticket 03 checked the continuation PDE with European Crank-Nicolson. Ticket 04 keeps
that same continuation machinery but adds the missing American ingredient: the payoff obstacle.

For a beginner, the key methodological shift is this:
\begin{itemize}
  \item European validation asks whether the finite-difference PDE machinery can reproduce a
        closed-form European benchmark.
  \item American pricing asks a constrained question: at each grid point, is it better to continue
        or to exercise immediately?
  \item The finite-difference version of that constrained question is an LCP, and PSOR is the
        first reusable numerical tool in this project for solving that LCP.
\end{itemize}

The formulation note describes the American problem as an obstacle problem. The validation plan
also emphasizes that later tickets should check obstacle violations and complementarity behavior.
Ticket 04 does not yet perform those full validation diagnostics, but it implements the solver core
that makes those diagnostics meaningful.

\section{Why American Options Become a Variational Inequality and LCP}

A European option can be exercised only at maturity. Between today and maturity, its value follows
the continuation PDE. In time-to-maturity notation, Ticket 03 solved
\begin{equation}
  U_{\tau} = L U,
\end{equation}
with terminal condition \(U(S,0)=\Phi(S)\).

An American option is different because the holder can exercise early. At every time and spot
point, the holder compares two choices:
\begin{itemize}
  \item continue holding the option;
  \item exercise immediately and receive payoff \(\Phi(S)\).
\end{itemize}

The value must therefore satisfy
\begin{equation}
  U(S,\tau) \geq \Phi(S).
  \label{eq:obstacle}
\end{equation}

In the continuation region, where holding is better than exercising, the PDE still applies. In the
exercise region, the value sits on the payoff. This is why the continuous problem is a
variational inequality rather than a plain PDE.

After finite-difference discretization and a Crank-Nicolson time step, the problem becomes a
linear complementarity problem, or LCP:
\begin{align}
  A u &\geq b, \\
  u &\geq \Phi, \\
  (u-\Phi)^T(Au-b) &= 0.
  \label{eq:lcp}
\end{align}

Here \(u\) is the vector of new interior-node values, \(A\) is the Crank-Nicolson left-hand
matrix, \(b\) is the adjusted right-hand side, and \(\Phi\) is the payoff vector on interior
nodes.

The three LCP lines each have a plain-language interpretation:
\begin{itemize}
  \item \(A u \geq b\): the discrete continuation equation is not allowed to point to a value
        that would make immediate exercise preferable.
  \item \(u \geq \Phi\): the option value must stay at or above the exercise payoff.
  \item \((u-\Phi)^T(Au-b)=0\): at each node, at least one of the two gaps is active. Either the
        node is above payoff and behaves like a continuation node, or it sits on payoff and the
        exercise constraint is binding.
\end{itemize}

A useful beginner picture is to think of two nonnegative gaps:
\begin{equation}
  \text{value gap}_i = u_i-\Phi_i,
  \qquad
  \text{equation gap}_i = (Au-b)_i.
\end{equation}
Complementarity says that their product should be zero at a properly solved node:
\begin{equation}
  \text{value gap}_i \times \text{equation gap}_i = 0.
\end{equation}
This does not mean both gaps must be zero. It means the node should not be simultaneously above
payoff and away from the continuation equation.

\section{Payoff Obstacle in Plain Language}

The payoff obstacle says that an American option cannot be worth less than what the holder can get
by exercising immediately. For a put,
\begin{equation}
  \Phi_{\mathrm{put}}(S)=\max(K-S,0).
\end{equation}
For a call,
\begin{equation}
  \Phi_{\mathrm{call}}(S)=\max(S-K,0).
\end{equation}

If a numerical solver reports \(U<\Phi\), the result is financially invalid. The holder would
simply exercise and receive \(\Phi\), so a lower computed value cannot be correct. In code,
Ticket 04 enforces this by projecting each PSOR update:
\begin{equation}
  u_i \leftarrow \max(\Phi_i,\; \text{SOR candidate}_i).
\end{equation}

This projection is the key difference between the European solver and the American solver. A
European solver can accept the linear-system answer because early exercise is not allowed. An
American solver must check the linear-system answer against the exercise payoff at every update.

Projection is local and mechanical, but it has a strong financial meaning. It says: after the row
calculation proposes a value, the algorithm asks whether an investor could do better by exercising
right now. If yes, the node is pulled back up to payoff. If no, the continuation value is allowed
to remain above payoff.

Projection alone is not a full validation proof. It protects against direct obstacle violation,
but later tickets still need to examine accuracy, complementarity residuals, and financial
benchmarks.

\section{Continuation Value Versus Immediate Exercise Value}

The continuation value is the value of keeping the option alive. The immediate exercise value is
the payoff. American pricing asks which is better at every grid point.

If continuation is better, then
\begin{equation}
  u_i > \Phi_i
\end{equation}
and the finite-difference equation should hold closely:
\begin{equation}
  (Au-b)_i \approx 0.
\end{equation}

If immediate exercise is better, then
\begin{equation}
  u_i = \Phi_i
\end{equation}
and the PDE equation is not forced in the same way. The complementarity condition in
Equation~\ref{eq:lcp} expresses this either-or logic. Each point should either behave like a
continuation point or sit on the obstacle.

\begin{longtable}{L{0.23\textwidth}L{0.31\textwidth}L{0.34\textwidth}}
\toprule
Region & Mathematical signal & Plain-language meaning \\
\midrule
\endhead
Continuation region &
\(u_i>\Phi_i\) and \((Au-b)_i\approx 0\) &
Holding the option is better than exercising. The node behaves like the European continuation
equation. \\
Exercise region &
\(u_i=\Phi_i\) and \((Au-b)_i\geq 0\) &
Immediate exercise is binding. The value is pinned to payoff, so the continuation equation is not
the active condition. \\
Near the boundary &
\(u_i-\Phi_i\) is small and numerical tolerances matter &
The grid point is close to the exercise frontier. Later tickets should interpret diagnostics with
tolerance rather than exact equality. \\
\bottomrule
\end{longtable}

For an American put, early exercise can become attractive when the spot price is low enough. For
an American call with dividends, early exercise can become structurally relevant because receiving
stock earlier may matter when dividends are present. Ticket 04 supports both paths structurally,
but it does not yet claim full validation for either case.

\section{PSOR and Projection in Plain Language}

Successive over-relaxation, or SOR, is an iterative way to solve a linear system. It updates one
row at a time and uses the newest values as soon as they are available. The relaxation parameter
\(\omega\) controls how strongly the solver moves toward the row solution:
\begin{equation}
  u_i^{\mathrm{SOR}} =
  u_i^{\mathrm{old}}
  + \omega\left(u_i^{\mathrm{GS}} - u_i^{\mathrm{old}}\right),
\end{equation}
where \(u_i^{\mathrm{GS}}\) is the ordinary Gauss-Seidel row candidate.

Projected SOR adds one extra step:
\begin{equation}
  u_i^{\mathrm{new}} =
  \max\left(\Phi_i,\; u_i^{\mathrm{SOR}}\right).
\end{equation}

In beginner language, each row asks: ``What value would the linear equation suggest?'' Then the
projection asks: ``Is that value below payoff?'' If it is below payoff, the update is pushed back
to payoff.

\begin{longtable}{L{0.24\textwidth}L{0.31\textwidth}L{0.33\textwidth}}
\toprule
Component & Role & Beginner interpretation \\
\midrule
\endhead
\(A\) & CN left-hand matrix & The continuation equation to solve at the new time level. \\
\(b\) & Adjusted right-hand side & Old value information plus boundary adjustments. \\
\(\Phi\) & Payoff vector & The minimum allowed American option value. \\
\(\omega\) & Relaxation parameter & How aggressively the row update moves. \\
Projection & \(u_i=\max(\Phi_i,\cdot)\) & The safeguard that enforces the obstacle. \\
Convergence metadata & Iterations and final update & Evidence about whether PSOR actually settled. \\
\bottomrule
\end{longtable}

\subsection{Step-by-Step PSOR Algorithm Walkthrough}

The projected SOR solve used in Ticket 04 is easier to understand as a repeated sweep through the
interior grid:
\begin{enumerate}
  \item Start from an initial guess. In the time-marching solver, the previous time-level values
        are a natural starting point.
  \item Sweep through the LCP rows from left to right.
  \item For the current row \(i\), compute the Gauss-Seidel row candidate. The left-neighbour
        value \(u_{i-1}\) has already been updated in the current sweep, so the newest available
        left value is used. The right-neighbour value \(u_{i+1}\) has not yet been visited in the
        current sweep, so its current stored value is used.
  \item Apply SOR relaxation. If \(\omega=1\), this is ordinary Gauss-Seidel. If
        \(1<\omega<2\), the update moves more aggressively toward the row candidate.
  \item Project above payoff:
        \[
          u_i \leftarrow \max(\Phi_i,u_i^{\mathrm{SOR}}).
        \]
        This is the American-option step.
  \item Track the largest absolute change made during the sweep.
  \item Stop when the largest update is less than the requested tolerance.
  \item If the maximum number of sweeps is reached first, report non-convergence instead of
        silently treating the answer as reliable.
\end{enumerate}

\begin{longtable}{L{0.20\textwidth}L{0.34\textwidth}L{0.34\textwidth}}
\toprule
PSOR stage & Numerical action & Beginner debugging question \\
\midrule
\endhead
Initial guess & Copy the supplied initial vector or start from payoff. & Is the starting vector the
right length and above payoff when expected? \\
Row candidate & Solve one tridiagonal row using latest left and current right neighbour values. &
Are the neighbour indices aligned with the matrix rows? \\
Relaxation & Blend old value and Gauss-Seidel candidate using \(\omega\). & Is
\(0<\omega<2\)? \\
Projection & Replace the candidate by \(\max(\Phi_i,\cdot)\). & Can any node fall below payoff
after the update? \\
Stopping check & Measure the maximum absolute row update. & Is non-convergence reported if the
tolerance is not reached? \\
\bottomrule
\end{longtable}

\section{How Ticket 04 Extends Ticket 03 European CN}

Ticket 03 used Crank-Nicolson in the European continuation case:
\begin{equation}
  \left(I-\frac{\Delta\tau}{2}L_h\right)U^{n+1}
  =
  \left(I+\frac{\Delta\tau}{2}L_h\right)U^n.
\end{equation}

Ticket 04 keeps the same matrix and boundary logic. It still builds
\begin{equation}
  A = I-\frac{\Delta\tau}{2}L_h
\end{equation}
and
\begin{equation}
  b^n = \left(I+\frac{\Delta\tau}{2}L_h\right)U^n
\end{equation}
with new-time boundary adjustments included in \(b^n\).

The change is the solve. Ticket 03 solved the unconstrained European linear system. Ticket 04
solves the constrained American LCP with PSOR:
\begin{align}
  A U^{n+1} &\geq b^n, \\
  U^{n+1} &\geq \Phi.
\end{align}

The Ticket 04 solver stores the full value grid over \((\tau,S)\) so later tickets can compute
obstacle diagnostics, continuation premiums, validation tables, and boundary extraction. Ticket
04 itself does not compute those later diagnostics.

\begin{longtable}{L{0.22\textwidth}L{0.34\textwidth}L{0.32\textwidth}}
\toprule
One CN/PSOR time-step stage & What the code is doing & Why it matters \\
\midrule
\endhead
Old value grid & Start with \(U^n\), the option values already computed at the previous
time-to-maturity level. & This is the known information for the next step. \\
Old boundary values & Use boundary formulas at \(\tau_n\) when building the right-hand side. &
Boundary nodes affect nearby interior nodes through the spatial operator. \\
Right-hand side construction & Form \((I+\frac{\Delta\tau}{2}L_h)U^n\) on the interior nodes. &
This carries old-time continuation information forward. \\
New boundary adjustment & Move known boundary contributions from \(A U^{n+1}\) into the
right-hand side. & The LCP solve handles only interior unknowns, so boundary terms must be
accounted for explicitly. \\
LCP solve by PSOR & Solve \(A u \geq b\), \(u\geq\Phi\) for the new interior vector. & This is
where American exercise enters the time step. \\
Projection against payoff & Each row update is clipped upward to payoff. & This enforces the
obstacle locally. \\
Metadata storage & Store convergence status, iteration count, tolerance, \(\omega\), and final
update. & Later validation can detect slow or failed PSOR steps. \\
\bottomrule
\end{longtable}

\section{LCP Matrix Form and Update Logic}

For a tridiagonal row,
\begin{equation}
  \ell_i u_{i-1} + d_i u_i + r_i u_{i+1} = b_i,
\end{equation}
where \(\ell_i\), \(d_i\), and \(r_i\) are the lower, diagonal, and upper matrix entries.

The Gauss-Seidel row candidate is
\begin{equation}
  u_i^{\mathrm{GS}} =
  \frac{
    b_i - \ell_i u_{i-1}^{\mathrm{new}} - r_i u_{i+1}^{\mathrm{old}}
  }{d_i}.
\end{equation}

Ticket 04 then applies relaxation and projection:
\begin{align}
  u_i^{\mathrm{SOR}} &=
  u_i^{\mathrm{old}} + \omega(u_i^{\mathrm{GS}}-u_i^{\mathrm{old}}), \\
  u_i^{\mathrm{new}} &=
  \max(\Phi_i,u_i^{\mathrm{SOR}}).
\end{align}

The implementation uses the maximum absolute update as a convergence proxy:
\begin{equation}
  \max_i |u_i^{\mathrm{new}}-u_i^{\mathrm{old}}|.
\end{equation}
When this quantity is below the requested tolerance, the PSOR solve reports convergence. If the
maximum iteration count is reached first, non-convergence is reported explicitly.

The words ``old'' and ``new'' are local to a PSOR sweep. During a single sweep, nodes already
visited on the left use their newest values, while nodes not yet visited on the right still use
their current stored values. This mixed use of newest and current values is not a bug; it is the
Gauss-Seidel idea that makes SOR different from a fully simultaneous Jacobi update.

\section{Implementation Design Decisions}

\subsection{One Focused Module}

Ticket 04 creates one source module, \file{src/american_risk_surfaces/solvers/cn_psor.py}. It
contains only the PSOR/LCP core and American CN/PSOR time-marching core. It does not create
diagnostic, boundary-extraction, Greek, experiment, dataset, stress-map, or neural-network modules.

\subsection{Interior Nodes Only}

The PSOR solver works on interior nodes only. Boundary values are imposed outside the LCP solve,
consistent with Tickets 02 and 03. This keeps the distinction clear: interior nodes are solved,
boundary nodes are assigned by boundary formulas.

\subsection{Metadata Is Part of the Result}

PSOR can fail to converge if the tolerance, relaxation parameter, grid, or maximum iteration count
is poor. Ticket 04 therefore returns convergence metadata rather than hiding this risk.

\begin{longtable}{L{0.24\textwidth}L{0.62\textwidth}}
\toprule
Metadata field & Meaning \\
\midrule
\endhead
\code{converged} & Whether the final update fell below tolerance before \code{max\_iter}. \\
\code{iterations} & Number of PSOR sweeps used. \\
\code{final\_update} & Maximum absolute update in the final sweep. \\
\code{tolerance} & Requested stopping tolerance. \\
\code{omega} & Relaxation parameter used in the solve. \\
\code{max\_iter} & Maximum number of sweeps allowed. \\
\bottomrule
\end{longtable}

\subsection{American Put and Dividend-Call Structure}

The time-marching solver supports American puts and American calls. It uses Ticket 01 payoff
functions and Ticket 02 American boundary helpers. The call path is structural at this stage: it
runs with dividend-aware call boundaries, but detailed no-dividend and dividend-paying call
validation belongs to later tickets.

\section{Code-Level Function Explanations}

\begin{longtable}{P{0.38\textwidth}L{0.54\textwidth}}
\toprule
Function or class & Purpose and important behavior \\
\midrule
\endhead
\code{PSORResult} &
Frozen dataclass containing one LCP solution and convergence metadata. \\
\code{AmericanCNPSORResult} &
Frozen dataclass containing option parameters, grids, payoff, full value grid, final values, per-step PSOR metadata, convergence status, and maximum obstacle violation. \\
\code{psor\_lcp\_solve} &
Solves a tridiagonal LCP on interior nodes using projected SOR. It validates array shapes, \(\omega\), tolerance, and max iterations. \\
\code{american\_crank\_nicolson\_psor\_price} &
Builds the grid and CN matrices, forms each right-hand side, solves each American time step with PSOR, and stores the value grid. \\
\bottomrule
\end{longtable}

The public API is controlled through \code{\_\_all\_\_}. Boundary helpers from Ticket 02 are used
internally but not exported as Ticket 04 public API.

\section{Testing Strategy}

\begin{longtable}{L{0.25\textwidth}L{0.31\textwidth}L{0.32\textwidth}}
\toprule
Test category & What it proves & What it does not prove \\
\midrule
\endhead
Toy LCP projection &
Projected SOR pushes values back above payoff in a known small system. &
Full option-pricing accuracy. \\
American put smoke test &
The basic put time-marching path preserves \(U\geq\Phi\) in a small case. &
Ticket 05 American put validation. \\
Metadata test &
Convergence fields are exposed and populated. &
That every parameter choice converges quickly. \\
Non-convergence test &
Failure to converge is reported rather than hidden. &
How to tune \(\omega\) for production grids. \\
\(T=0\) test &
The solver returns payoff at maturity without PSOR steps. &
Positive-maturity American behavior. \\
Dividend call structure test &
The call path runs with call payoff and call-style lower boundary. &
No-dividend or dividend-paying call validation. \\
Invalid input tests &
Bad option type, grid inputs, \(\omega\), tolerance, and LCP shapes raise \code{ValueError}. &
Every possible production input policy. \\
Scope test &
The public API does not add boundary extraction, Greeks, stress, dataset, or neural surfaces. &
Future tickets that intentionally add those features. \\
\bottomrule
\end{longtable}

\section{Test Results and Interpretation}

The exact test command was:
\begin{lstlisting}[language=bash]
PYTHONPYCACHEPREFIX=/private/tmp/codex_pycache PYTHONPATH=src python3 -m unittest discover -s tests
\end{lstlisting}

Observed result after implementing Ticket 04:
\begin{lstlisting}
Ran 40 tests in 0.030s

OK
\end{lstlisting}

The compile command also completed successfully:
\begin{lstlisting}[language=bash]
PYTHONPYCACHEPREFIX=/private/tmp/codex_pycache PYTHONPATH=src python3 -m compileall -q src tests experiments
\end{lstlisting}

For the American put smoke case used by the tests, the solver reported convergence across 40 time
steps, maximum obstacle violation \(0\), and no hidden non-convergence. This is a useful local
correctness check for the PSOR core. It is not a full validation study.

\begin{longtable}{L{0.30\textwidth}L{0.26\textwidth}L{0.32\textwidth}}
\toprule
Observed outcome & Recorded result & Interpretation \\
\midrule
\endhead
Total unittest count &
40 tests passed &
The Ticket 04 tests ran together with the earlier Ticket 01--03 tests. \\
American put smoke case &
Converged across 40 time steps &
The basic American put CN/PSOR path completed every step in the small test case. \\
Maximum obstacle violation &
\(0\) &
No tested grid value fell below payoff in the American put smoke case. \\
Non-convergence case &
Reported explicitly &
A deliberately constrained PSOR run returned \code{converged=False} rather than hiding the
failure. \\
\(T=0\) behavior &
Returned payoff &
At maturity, the American option value equals immediate exercise payoff and no time marching is
needed. \\
\bottomrule
\end{longtable}

These outcomes are important, but narrow. They show that the implementation respects the main
mechanical requirements of Ticket 04. They do not establish production accuracy for American puts,
do not compare against external reference values, and do not prove PSOR convergence for all
parameter choices.

\section{Implementation Pitfalls and Debugging Notes}

Ticket 04 recorded one expected test-driven red step and one workflow-level import issue. Neither
was a mathematical failure of the PSOR algorithm, but both are useful for beginner researchers to
understand because they show how implementation checks can fail for ordinary engineering reasons.

\begin{longtable}{L{0.22\textwidth}L{0.24\textwidth}L{0.21\textwidth}L{0.19\textwidth}}
\toprule
Recorded item or pitfall & Symptom & Fix or guard & Beginner lesson \\
\midrule
\endhead
Expected TDD red run before implementation &
\code{ModuleNotFoundError} for the missing \code{cn\_psor} module. &
Create the Ticket 04 module and public functions, then rerun the same tests. &
A red test is useful when it proves the test targets the missing API. \\
Ad-hoc metadata probe without \code{PYTHONPATH=src} &
Package import failed even though the source tree existed. &
Rerun using the project test convention \code{PYTHONPATH=src}. &
Import failures can be environment issues, not code logic issues. \\
Mixing PSOR neighbour timing incorrectly &
Updates may converge slowly or solve the wrong row equation. &
Use newest left-neighbour values and current right-neighbour values during each sweep. &
Gauss-Seidel/SOR is sequential, not simultaneous. \\
Wrong boundary adjustment sign &
Interior values near boundaries can be distorted. &
Reuse the Ticket 03 CN boundary-adjustment pattern carefully. &
Boundary terms are known values, but they still affect the interior system. \\
Forgetting projection above payoff &
The computed American value can fall below \(\Phi\). &
Apply \(u_i=\max(\Phi_i,\text{candidate})\) at every row update. &
Projection is the American-option part of PSOR. \\
Hiding non-convergence &
The solver may look successful after reaching \code{max\_iter}. &
Return \code{converged=False}, iteration count, and final update. &
Numerical algorithms need honest status metadata. \\
Using put payoff in the call path &
The dividend-call structure would run with the wrong exercise value. &
Test the call path separately with call payoff and call-style boundaries. &
American puts and calls share machinery but not payoff logic. \\
Treating Ticket 04 smoke tests as full validation &
Small tests could be overinterpreted as solver accuracy evidence. &
State clearly that full American put validation belongs to Ticket 05. &
Passing smoke tests is necessary, but not enough for model credibility. \\
Overclaiming from a small smoke test &
The report could imply the whole American solver is validated. &
Keep limitations explicit and avoid claiming benchmark accuracy. &
Good computational finance reports separate implementation checks from validation studies. \\
\bottomrule
\end{longtable}

Future coding and testing reports should include an ``Implementation Pitfalls and Debugging
Notes'' section when useful, especially if failed tests, environment mistakes, numerical bugs, or
fixes occurred during implementation. This makes the reports more valuable as study material
instead of only recording the final clean state.

\section{Limitations}

\begin{longtable}{L{0.34\textwidth}L{0.54\textwidth}}
\toprule
Not included in Ticket 04 & Reason \\
\midrule
\endhead
Full American put validation & Belongs to Ticket 05. \\
No-dividend American call validation & Belongs to a later call-validation ticket. \\
Dividend-paying American call validation & Belongs to a later dividend-call ticket. \\
Boundary extraction & Requires continuation-premium logic planned for Ticket 09. \\
Greeks & Delta and Gamma diagnostics are later work. \\
Stress maps, datasets, neural networks & These depend on a fully validated solver. \\
Figures & Not essential for this core-code ticket. \\
Formal PSOR convergence proof & Background theory, not part of this implementation ticket. \\
\bottomrule
\end{longtable}

The current implementation uses a final-update stopping proxy. Later diagnostic tickets may add
more detailed complementarity residual calculations and reporting.

\section{Lessons Learned / Study Notes}

The first lesson is that American option pricing is constrained pricing. The value is not allowed
to fall below payoff because the holder can exercise immediately.

The second lesson is that the LCP is the discrete version of the exercise decision. A node is
either a continuation node, where the equation is active, or an exercise node, where the payoff
obstacle is active.

The third lesson is that PSOR is not just a linear solver with a different name. Projection is the
important American-option step. Without projection, the method would only be a European
continuation solve.

The fourth lesson is that convergence metadata matters. If PSOR stops too early or reaches the
maximum iteration count, the solver output should say so directly.

The fifth lesson is to keep validation layered. Ticket 04 checks the reusable PSOR mechanism and
basic American time-marching structure. It does not replace the detailed American put validation
that comes next.

\section{Link to Ticket 05 American Put Validation}

Ticket 05 should use the Ticket 04 solver to run a detailed American put validation study. That
next ticket should add validation tables, selected price checks, obstacle summaries, and more
financial interpretation for the put case. It should still avoid boundary extraction and Greeks
unless the approved roadmap explicitly moves to those later tickets.

The handoff from Ticket 04 to Ticket 05 is:
\begin{itemize}
  \item Ticket 04 provides the American CN/PSOR solver and metadata.
  \item Ticket 05 should evaluate whether the American put results are numerically and financially
        credible.
\end{itemize}

\section{Reproducibility Record}

\subsection{Files Created}

\begin{longtable}{P{0.52\textwidth}L{0.40\textwidth}}
\toprule
Path & Role in Ticket 04 \\
\midrule
\endhead
\file{src/american_risk_surfaces/solvers/cn_psor.py} &
PSOR/LCP core and basic American CN/PSOR time-marching solver. \\
\file{tests/test_psor.py} &
Focused \code{unittest} coverage for projected SOR, American smoke behavior, metadata, non-convergence, invalid inputs, and scope. \\
\file{reports/03_solver/tickets/ticket_04_psor_lcp_core.tex} &
Formal LaTeX study-report source. \\
\file{reports/03_solver/tickets/ticket_04_psor_lcp_core.pdf} &
Formal PDF compiled from the LaTeX source. \\
\bottomrule
\end{longtable}

\subsection{Commands}

Test command:
\begin{lstlisting}[language=bash]
PYTHONPYCACHEPREFIX=/private/tmp/codex_pycache PYTHONPATH=src python3 -m unittest discover -s tests
\end{lstlisting}

Python compile command:
\begin{lstlisting}[language=bash]
PYTHONPYCACHEPREFIX=/private/tmp/codex_pycache PYTHONPATH=src python3 -m compileall -q src tests experiments
\end{lstlisting}

LaTeX compile command:
\begin{lstlisting}[language=bash]
HOME=/private/tmp FC_CACHEDIR=/private/tmp/fontconfig-cache latexmk -xelatex -interaction=nonstopmode -halt-on-error -outdir=/private/tmp/ticket04_latex reports/03_solver/tickets/ticket_04_psor_lcp_core.tex
\end{lstlisting}

PDF verification uses \code{pdfinfo} and page rendering with \code{pdftoppm}. Raw Git status is
intentionally not included in this formal PDF report.

\bibliographystyle{plainnat}
\bibliography{reports/03_solver/references}

\end{document}
